% !TEX TS-program = lualatex
Рассмотрим более детально строение модуля игры (\refimage{game-core-struct}):

\image[width=0.9\textwidth]{Строение модуля игры}{game-core-struct}{game-core-struct}

Класс BackgammonEnv наследуется от интерфейса gym.Env и реализует необходимые методы (step, reset, render и др.), обеспечивая совместимость с алгоритмами обучения с подкреплением. Он управляет взаимодействием между агентом и логикой игры, обрабатывая действия, полученные от агента, и возвращая соответствующие наблюдения, награды и признаки окончания эпизода. Внутри класса используется экземпляр класса Backgammon, который инкапсулирует саму игровую механику.

Класс Backgammon представляет собой ядро всей игровой логики. Он отвечает за генерацию допустимых ходов, выполнение действий, определение победителя и сохранение и восстановление состояния игры. Важно отметить наличие вспомогательных методов, таких как get\_normal\_plays, get\_bear\_off\_plays, и get\_bar\_plays, которые обрабатывают различные фазы игры (выход с бара, снятие фишек и обычные ходы).

Состояние игры описывается с помощью структуры BackgammonState, реализованной на основе namedtuple. Это позволяет эффективно сохранять и восстанавливать игру на любом этапе.

Класс Constants содержит набор фиксированных значений, таких как идентификаторы игроков (WHITE, BLACK), количество пунктов на доске и специальные зоны (BAR, OFF), что обеспечивает читаемость и удобство поддержки кода.

Для взаимодействия с окружением может быть использован внешний агент, как, например, класс RandomAgent, реализующий простую стратегию выбора действий. Он демонстрирует, каким образом можно интегрировать пользовательские агенты в процесс обучения или тестирования.
